\section{Exact Greedy Selection}
\label{app:greedy}

\begin{proposition}[The rounds are exact greedy]
\label{prop:greedy-exact}
Fix a candidate set $\mathcal{C}'\subseteq\mathcal{C}$ and a smoothing constant $\gamma$, and let $w_j=\lambda^{\,j-1}$.
\begin{enumerate}[label=(\roman*),leftmargin=*]
\item The score $\HD_\gamma(p,A)\cdot\lambda^{Counts[A]}$ that line~4 of Algorithm~\ref{alg:topk} assigns a remaining candidate is exactly its marginal gain in $F_w$, so selecting the highest-scoring candidate is one step of greedy on $F_w$.
\item For any non-increasing $w$, selecting $E\subseteq\mathcal{C}'$ with $|E|=k$ maximising $F_w(E)$ is solvable in polynomial time, and greedy returns an optimal solution.
\end{enumerate}
\end{proposition}

\begin{proof}
(i) In a given round all candidates carrying an attribute $A$ share the factor $\lambda^{Counts[A]}$, so whenever a round selects from that group it selects the member of largest $\HD_\gamma$; after $c$ selections from $A$ the members taken are therefore its $c$ highest-scoring ones.
Every remaining candidate of $A$ thus scores no higher than any already taken, so adding it places it at rank $c+1$ of its group and leaves the ranks of the existing members unchanged.
The group term of Eq.~\eqref{eq:set-objective} gains exactly $w_{c+1}h=\lambda^{c}h$, which is the score of line~4, and no other group term changes.

(ii) By Definition~\ref{def:setobj}, $F_w(E)=\sum_A g_A(E)$ where $g_A$ depends only on the members of $E$ whose attribute is $A$.
There is therefore no interaction between groups, and a solution is determined by how many candidates $n_A\ge0$ it takes from each group, subject to $\sum_A n_A=k$, together with which ones.
For a fixed $n_A$, the best choice within group $A$ is its $n_A$ highest-scoring candidates, since $w$ is non-increasing and the scores are ranked (rearrangement inequality).
Hence the value of taking $g$ candidates from group $A$ is
\[
f_A(g)=\sum_{j=1}^{g} w_j\, h^A_{(j)},
\]
whose successive marginal gains are $f_A(g)-f_A(g-1)=w_g\,h^A_{(g)}$.
These are non-increasing in $g$, because $w_g$ is non-increasing by assumption and $h^A_{(g)}$ is non-increasing by construction; so each $f_A$ is concave.

Maximising a separable sum of concave functions subject to a cardinality constraint is exactly the problem solved by the greedy rule of repeatedly taking the largest remaining marginal gain, which is optimal for separable concave resource allocation.
Since each marginal gain is computable in constant time after sorting each group, the total cost is $O(|\mathcal{C}'|\log|\mathcal{C}'|)$ for $|\mathcal{C}'|$ candidates.
With (i), the $k$ rounds of Algorithm~\ref{alg:topk} run exactly this procedure.
\end{proof}



\begin{example}[Greedy and the optimum agree]
\label{ex:greedy}
Let $\lambda=\tfrac12$ and $k=3$, and let five candidates carry the scores $9$ and $4$ on one attribute and $7$, $6$ and $2$ on another.
The marginal gains available are $9$ and $\tfrac12\cdot4=2$ in the first group, and $7$, $\tfrac12\cdot6=3$ and $\tfrac14\cdot2=0.5$ in the second.
Greedy takes $9$, then $7$, then $3$, for a total of $19$.
These are the three largest gains in the list, so no set of three candidates can do better, and enumerating all ten subsets confirms that $19$ is the optimum.
Taking the two best scores of the first group instead, $9$ and $4$, together with $7$, gives $18$: the second copy of an attribute is discounted, which is what makes the greedy order follow the discounted gains rather than the raw scores.
\end{example}
